\section*{Capítulo \ref{ch:b1:elementary-steps} --- Mecanismos de reação: etapas elementares e aproximações}

\begin{solution}{exo:b1:elementary-steps:1}
(a) Unimolecular: pode ser elementar. (b) Bimolecular: pode ser elementar.
(c) Três moléculas e quatro ligações reorganizadas de uma vez: não é
elementar. (d) Termolecular, com o \ce{N2} levando embora a energia
liberada: uma \omterm{def:b1:elementary-steps:elementary-step}{etapa elementar} rara, mas verdadeira.
\end{solution}

\begin{solution}{exo:b1:elementary-steps:2}
$v = k[\mathrm{A}][\mathrm{B}]$; $v = k[\mathrm{A}]^2$; $v = k[\mathrm{A}]$.
Uma equação global é um balanço, não uma descrição de como as moléculas
se encontram; sua \omterm{def:b1:rate-laws:order}{lei de velocidade} precisa ser medida, ou deduzida de
um mecanismo.
\end{solution}

\begin{solution}{exo:b1:elementary-steps:3}
O intermediário é o vale em 30; os \omterm{def:b1:elementary-steps:energy-profile}{estados de transição} são os máximos
em 70 e 55. O \omterm{def:b1:elementary-steps:energy-profile}{estado de transição} mais alto é o primeiro: a primeira
etapa é a determinante. Ela é endotérmica (de 0 a 30).
\end{solution}

\begin{solution}{exo:b1:elementary-steps:4}
Pelo mesmo fator $10^4$ (\cref{prop:b1:elementary-steps:catalysis}); a
constante de equilíbrio não muda; o equilíbrio é atingido cerca de
$10^4$ vezes mais cedo.
\end{solution}

\begin{solution}{exo:b1:elementary-steps:5}
$t_{\max} = \ln(0.30/0.10)/(0.30 - 0.10) = \qty{5.49}{s}$. Então
$[\mathrm{B}]_{\max} = a_0\frac{0.10}{0.20}(\eu^{-0.549} - \eu^{-1.648}) =
0.5\,a_0(0.577 - 0.192) = 0.192\,a_0$.
\end{solution}

\begin{solution}{exo:b1:elementary-steps:6}
$\dd[\mathrm{I}]/\dd t = k_1[\mathrm{A}] - (k_{-1} + k_2)[\mathrm{I}] = 0$
dá $[\mathrm{I}] = k_1[\mathrm{A}]/(k_{-1} + k_2)$ e $v = k_2[\mathrm{I}]
= \frac{k_1k_2}{k_{-1}+k_2}[\mathrm{A}]$. Se $k_2 \gg k_{-1}$, $v =
k_1[\mathrm{A}]$: a primeira etapa é a determinante. Se $k_2 \ll k_{-1}$,
$v = (k_1/k_{-1})k_2[\mathrm{A}]$: um pré-equilíbrio seguido de uma etapa
lenta.
\end{solution}

\begin{solution}{exo:b1:elementary-steps:7}
Soma: \ce{2H2O2 -> 2H2O + O2}. \omterm{def:b1:elementary-steps:catalyst}{Catalisador}: \ce{I-} (consumido e depois
regenerado); intermediário: \ce{IO-}. A primeira etapa, lenta, fixa a
velocidade: $v = k_1[\ce{H2O2}][\ce{I-}]$.
\end{solution}

\begin{solution}{exo:b1:elementary-steps:8}
A etapa é fortemente ascendente: pelo postulado de Hammond, seu \omterm{def:b1:elementary-steps:energy-profile}{estado
de transição} se assemelha ao carbocátion. Um substituinte que abaixa a
energia do carbocátion abaixa quase tanto a do \omterm{def:b1:elementary-steps:energy-profile}{estado de transição}, de
modo que a barreira diminui e a etapa fica mais rápida.
\end{solution}

\begin{solution}{exo:b1:elementary-steps:9}
$k_{\mathrm B}/k_{\mathrm C} = \exp(20000/RT)$: \num{3.0e3} a \qty{300}{K},
$\exp(4.01) = 55$ a \qty{600}{K}. O aquecimento reduz a preferência
cinética por B e, ao tornar reversível a formação de B, permite que o
sistema chegue ao C, mais estável.
\end{solution}

\begin{solution}{exo:b1:elementary-steps:10}
$k_1[\mathrm{A}][\mathrm{M}] = (k_{-1}[\mathrm{M}] + k_2)[\mathrm{A}^*]$,
logo $v = k_2[\mathrm{A}^*] = \dfrac{k_1k_2[\mathrm{A}][\mathrm{M}]}{k_{-1}[\mathrm{M}]
+ k_2}$. A alta pressão ($k_{-1}[\mathrm{M}] \gg k_2$), $v =
(k_1k_2/k_{-1})[\mathrm{A}]$: ordem 1. A baixa pressão, $v =
k_1[\mathrm{A}][\mathrm{M}]$: ordem 2; a ativação por colisão torna-se a
etapa lenta.
\end{solution}

\begin{solution}{exo:b1:elementary-steps:11}
A demonstração é a do teorema. Para $k_1 = k_2 = k$:
$\frac{\dd}{\dd t}([\mathrm{B}]\eu^{kt}) = ka_0$, logo $[\mathrm{B}]\eu^{kt} =
ka_0t$ e $[\mathrm{B}] = a_0kt\,\eu^{-kt}$, máxima em $t = 1/k$.
\end{solution}

\begin{solution}{exo:b1:elementary-steps:12}
$\mathrm{A} + \mathrm{B} \rightleftharpoons \mathrm{I} + \mathrm{C}$ (rápida,
constante $K_1$), depois $\mathrm{I} + \mathrm{B} \to \mathrm{P}$ (lenta,
$k_2$). A soma é a reação global; $[\mathrm{I}] =
K_1[\mathrm{A}][\mathrm{B}]/[\mathrm{C}]$ e $v = k_2[\mathrm{I}][\mathrm{B}]
= k_2K_1[\mathrm{A}][\mathrm{B}]^2/[\mathrm{C}]$: o subproduto, liberado
no pré-equilíbrio, desacelera a reação.
\end{solution}

\begin{solution}{pb:b1:elementary-steps:1}
\textbf{1.} Dobrar $[\ce{NO}]_0$ multiplica $v_0$ por 4: ordem 2; dobrar
$[\ce{O2}]_0$ a multiplica por 2: ordem 1.
\textbf{2.} $v = k[\ce{NO}]^2[\ce{O2}]$; $k = 1.0 \times 10^{-5}/((1.0
\times 10^{-3})^2 \times 1.0 \times 10^{-3}) =
\qty{1.0e4}{L^2.mol^{-2}.s^{-1}}$.
\textbf{3.} Uma colisão simultânea de três moléculas é rara, e uma única
etapa não poderia ficar mais lenta com o aquecimento (questão 5).
\textbf{4.} $E_a = R\ln(k_{400}/k_{300})/(1/300 - 1/400) = 8.314 \ln(0.1)
/(8.33 \times 10^{-4}) = \qty{-23}{kJ/mol}$.
\textbf{5.} Uma \omterm{def:b1:elementary-steps:elementary-step}{etapa elementar} ocorre por meio de colisões que
atravessam uma barreira; a fração delas capaz de fazê-lo,
$\eu^{-E_a/RT}$ com $E_a \ge
0$, só pode crescer com a temperatura.
\textbf{6.} $-\dd[\ce{NO}]/\dd t = 2v$, $-\dd[\ce{O2}]/\dd t = v$.
\textbf{7.} \ce{2NO -> N2O2} mais \ce{N2O2 + O2 -> 2NO2} dá
\ce{2NO + O2 -> 2NO2}; o intermediário é \ce{N2O2}.
\textbf{8.} Ambas bimoleculares (e o inverso da primeira, unimolecular).
\textbf{9.} $v = k_2[\ce{N2O2}][\ce{O2}]$.
\textbf{10.} $[\ce{N2O2}] = K_1[\ce{NO}]^2$ com $K_1 = k_1/k_{-1}$.
\textbf{11.} $v = k_2K_1[\ce{NO}]^2[\ce{O2}]$: a lei observada.
\textbf{12.} $k = k_1k_2/k_{-1}$.
\textbf{13.} $\dd[\ce{N2O2}]/\dd t = k_1[\ce{NO}]^2 - k_{-1}[\ce{N2O2}] -
k_2[\ce{N2O2}][\ce{O2}]$.
\textbf{14.} Igualando a zero: $[\ce{N2O2}] = k_1[\ce{NO}]^2/(k_{-1} +
k_2[\ce{O2}])$.
\textbf{15.} $v = k_2[\ce{N2O2}][\ce{O2}]$ dá a lei enunciada.
\textbf{16.} $k_{-1} \gg k_2[\ce{O2}]$: o dímero se desfaz com muito mais
frequência do que encontra o dioxigênio, de modo que a primeira etapa
permanece em equilíbrio.
\textbf{17.} Se $k_2[\ce{O2}] \gg k_{-1}$, $v = k_1[\ce{NO}]^2$: ordem 2
em \ce{NO} e 0 em \ce{O2}.
\textbf{18.} A ordem 1 medida em \ce{O2}: o primeiro limite.
\textbf{19.} $\ln k = \ln k_1 + \ln k_2 - \ln k_{-1}$.
\textbf{20.} Derivando em relação a $T$ e usando $\dd\ln k_i/\dd T
= E_i/RT^2$: $E_a = E_1 + E_2 - E_{-1}$.
\textbf{21.} O aquecimento acelera um pouco a etapa lenta ($E_2$ pequena),
mas desloca fortemente o pré-equilíbrio de volta para \ce{NO} ($E_{-1}$
grande): menos dímeros, e a velocidade líquida cai.
\textbf{22.} O dímero fica $E_1 - E_{-1} = \qty{-38}{kJ/mol}$ abaixo dos
reagentes, e o segundo \omterm{def:b1:elementary-steps:energy-profile}{estado de transição}, $E_2 = \qty{15}{kJ/mol}$
acima do dímero: \qty{23}{kJ/mol} abaixo dos reagentes. O ponto mais alto
do caminho é o primeiro \omterm{def:b1:elementary-steps:energy-profile}{estado de transição}, só \qty{2}{kJ/mol} acima.
\textbf{23.} $\exp\big(\frac{23000}{8.314}(\frac{1}{400} - \frac{1}{300})\big)
= 0.10$: dez vezes mais lenta a \qty{400}{K}.
\textbf{24.} $E_a = 2 - 40 + 15 = \textbf{\qty{-23}{kJ/mol}}$, o valor
medido na Parte I: o mecanismo explica tanto a \omterm{def:b1:rate-laws:order}{lei de velocidade} quanto
sua estranha dependência com a temperatura.
\end{solution}
